\section*{Methods}
\label{section: main paper data processing}
\doublespacing
\begin{comment}
\begin{center}
\begin{tabular}{|m{5em} | m{14em} | m{5em} | m{14em}|} 
 \hline
 Symbol & Description & Symbol & Description \\ [0.5ex] 
 \hline
 $U_{g \times s \times c}$ & Observed gene expression for a specific gene $g$ at spot $s$ for cell type $c$ & $ISM_{ls_1,ts_2}$ & Inherent Similarity Measure between ligand $l$ at spot $s_1$ and target $t$ at spot $s_2$ \\ 
 \hline
 $M$ & A vector with elements $\{m_g\}$ representing gene-specific scaling parameters & $ES_{ls_1,ts_2}$ & Edge Score between ligand $l$ at spot $s_1$ and target $t$ at spot $s_2$ \\
 \hline
 $W$ & A matrix with elements $\{w_{s,c}\}$ representing the number of cells of cell type $c$ in spots $s$ & $NS_{l,t}(s)$ & Neighborhood Score between ligand $l$ and target $t$ at spot $s$ \\
 \hline
 $G$ & $C \times G$ matrix of reference cell type signatures with elements $\{g_{c,g}\}$, which consist of $c = \{1, ... C\}$ gene expression profiles for $g = \{1, ... G\}$ genes & $r_{x,y} / r_{u,v,p,q}$ & pearson correlation between $x$ and $y$ $/$ average pearson correlation between $u$ and $p$ and $v$ and $q$ \\
 \hline
 $E_{g \times c \times s}$ & Expected gene expression for a specific gene $g$ at spot $s$ for cell type $c$ & $L_{X, Y}$ & Lees statistic bivariate spatial association measure for two variables $X$ and $Y$ \\
 \hline
 $ECS_{ls_1,ts_2}$ & Expression Correspondence Score between ligand $l$ at spot $s_1$ and target $t$ at spot $s_2$ & $avg(n, M)$ & average expression/abundance of $n$ across a set of spots $M$ \\
 \hline
\end{tabular}
\end{center}
\end{comment}

\subsection*{Overview of Renoir}
Using either paired cell type annotated scRNA-seq data and spatial transcriptomic (ST) data or single-cell resolution spatial transcriptomic data, and a set of curated ligand-target pairs, Renoir calculates spatially enriched neighbourhood scores for the curated set of ligand-target pairs for a given sample. The neighbourhood scores can subsequently be leveraged to track ligand-target activity across the geographic topology, as well as identify spatially enriched communication regions, cell type-specific interactions, and pathway activity.

\subsection*{Required input data}
\paragraph{Single-cell RNA-seq and Spatial transcriptomic data:} As input, Renoir requires the following - (1) a spatial transcriptomics dataset $X^{st} \in \mathbb{R}^{S \times G}$ with gene expression for $S$ spots and $G$ genes. Depending on the technology, the spots can contain a mixture of cells (e.g., 10X Visium) or single-cell (e.g., Nanostring CosMX). In case of cell-resolution ST data, cell type annotations of the spots to one of $C$ cell types is required. (2) Spatial coordinates for each spot, and (3) a corresponding cell type annotated single-cell gene expression dataset $X^{sc} \in \mathbb{R}^{N_c \times G}$ for $N_c$ cells and $G$ genes (required for ST datasets where each spot contains mixture of cells), where each cell mapped to one of $C$ cell types. Note that the single-cell gene expression dataset is optional when single-cell resolution spatial transciptomic dataset is available.

% a single-cell RNA-seq dataset $X^{sc}$ with $N_c$ cells by $G$ genes, a corresponding spatial transcriptomics dataset $X^{st}$ of $S$ spots by $G$ genes with coordinate positions of each spot and a mapping for each cell within the single cell dataset to one of $C$ cell type annotations. 
\paragraph{Acquiring curated ligand-target pairs:} We first curate a set of secreted ligands from NATMI's connectomeDB2020 \citep{Hou2020} database. For each ligand, we obtain a list of putative targets sorted according to the descending order of ligand–target regulatory potential scores as calculated by NicheNet's prior knowledge model \citep{Browaeys2020}. The number of targets per ligand is user-defined. For our analyses, we have used the top $100$ targets for each ligand. The ligand-target pair is retained for further analysis if both genes are among the highly variable genes in the spatial transcriptomic dataset. Thus, we obtain a set of ligand-target pairs, $\CS_{lt}$ specific to our spatial transcriptomic dataset. 

% \textcolor{blue}{ligand-receptor pairs from an existing compendium of available lists of ligand-receptor pairs in literature \citep{Armingol2021}. We then select a set of 

% The ligands are fed into  literature comprehensive prior model [reference] and the top $n \in \{10,100,500\}$ targets were selected for each ligand. The ligand-target pairs are further filtered by intersecting them with genes present within the sample to acquire a dataset specific set of ligand-target pairs.\\

\subsection*{Calculation of neighborhood activity score for a ligand-target pair for a spot}

\paragraph{Constructing neighborhoods:} We utilize the spatial transcriptomic data to construct a directed graph such that each spot within the dataset represents a node in the graph and there exists a directed edge between a node (spot) and every node within its neighborhood. A neighborhood for a spot is defined with respect to the distribution of spots within the spatial dataset considered. For data wherein the spots are equidistant from their immediate neighbors (e.g., 10X Visium and Spatial Transcriptomics), the neighborhood for a spot $s$ is defined to be a set of spots inclusive of $s$ such that every spot other than $s$ is at a ``one-hop" distance away from $s$. For data wherein the spots are not uniformly distributed across the sample (e.g., Slide-seq, StereoSeq, CosMx), the neighborhood for a spot $s$ is defined to be a set of spots (inclusive of $s$) which are at most $r$ unit distance away from $s$, where $r$ is user defined.

\paragraph{Formulating neighborhood activity scores:} On constructing the neighborhood graph, we quantify the potential interaction between a ligand $l$ and target $t$ within the neighborhood of a spot $s$ using a neighborhood activity score defined by,
\begin{equation}
    \CA_{l,t}(s) = ES_{ls,ts} + \Sigma_{j \in \CN_{s} \setminus s} (ES_{lj,ts} + ES_{ls,tj})
\end{equation}

where $\CN_s$ denotes the neighborhood of $s$ and $ES_{ls_1,ts_2}$ denotes the edge score (edge weight) of an edge directed from a spot $s_1$ expressing ligand $l$ to a neighboring spot $s_2$ expressing target $t$ and is defined by,
\\
$$ES_{ls_1,ts_2} = norm(ECS_{ls_1,ts_2}) \times norm(ISM_{ls_1,ts_2})$$
\\
where $norm(ECS_{ls_1,ts_2})$ and $norm(ISM_{ls_1,ts_2})$ are the min-max normalized expression correspondence score and inherent similarity measure for the ligand-target pair $l-t$ for the spots $s_1$ and $s_2$ respectively. The edge score quantifies the possible interaction between a ligand $l$ at spot $s_1$ and target $t$ at spot $s_2$ as the product of the inferred similarity between the expressions of $l$ and $t$ across all the cell types present at $s1$ and $s2$ respectively ($ECS_{ls_1,ts_2}$) and the inherent similarity (acquired from the single cell data) between the expressions of $l$ and $t$ amongst cell types present at $s_1$ and $s_2$ respectively.

\paragraph{Calculation of Expression Correspondence Score and Inherent Similarity Measure:} To calculate the Expression Correspondence Score and Inherent Similarity Measure between two neighboring spots $s_1$ and $s_2$ where a ligand $l$ is expressed at spot $s_1$ and target $t$ is expressed at spot $s_2$, we utilize the following formulations
\begin{equation}
    ECS_{ls_1,ts_2} = \frac{\Sigma_{c_1 \in C_1} p_{c_1}PEM_{l,c_1,s_1} + \Sigma_{c_2 \in C_2} p_{c_2}PEM_{t,c_2,s_2}}{2}
\end{equation}

\begin{equation}\label{e4}
    ISM_{ls_1,ts_2} = \frac{\Sigma_{c_1 \in C_1} -p_{c_1}log_{10}(\frac{MI(lc_1,tc_1)}{min(H_{lc_1},H_{tc_1})}) + \Sigma_{c_2 \in C_2} -p_{c_2}log_{10}(\frac{MI(lc_2,tc_2)}{min(H_{lc_2},H_{tc_2})})}{2}
\end{equation}

%\\
%or,\\
%\\
%$$ECS_{ls_1,ts_2} = \Sigma_{c_1 \in C_1, c_2 \in C_2} \frac{1}{|p_{c_1} PEM_{l,c_1,s_1} - p_{c_2} PEM_{t,c_2,s_2}|}$$

where, $C_1$ and $C_2$ denote the set of cell types present at $s_1$ and $s_2$ respectively, $p_{c_1}$ and $p_{c_2}$ denote the cell type proportions for cell types $c_1 \in C_1$ and $c_2 \in C_2$ and $c_2$ expresses at least one receptor of $l$ respectively. $c_2$ is considered to express a receptor if the receptor gene is expressed in more than $thresh$\% of cells of that cell type (computed based on scRNA-seq/single-cell ST data), where $thresh$ in our analysis was set to 5 following NicheNet. $ECS_{ls_1,ts_2}$ is defined as the average of cell type proportion-weighted PEM of ligand $l$ and target $t$ in spots $s_1$ and $s_2$ respectively. The PEM value for cell type specificity of gene $g$ in cell type $c$ for spot $s$ is given by 
\begin{equation}
    PEM_{g,c,s} = log_{10}(\frac{u_{scg}}{e_{scg}}),
\end{equation}
where, $u_{scg}$ and $e_{scg}$ denote the inferred and expected expression of gene $g$ specific to cell type $c$ at spot $s$ respectively. In \cref{e4}, $MI(lc_1,tc_2)$ denotes the mutual information between the expression of $l$ across cell type $c_1$ and the expression of $t$ across cell type $c_2$. $H(lc_1)$ denotes the Shannon entropy of the expression of ligand $l$ across cell type $c_1$ and $H(tc_2)$ denotes the Shannon entropy of the expression of target $t$ across cell type $c_2$. $MI(lc_1,tc_2)$, $H(lc_1)$, and $H(tc_2)$ have been estimated from the single-cell (either scRNA-seq or single-cell ST) data using the Miller-Madow correction.

\paragraph{Quantifying inferred and expected cell type-specific gene expression:} In order to quantify the cell type-specific gene expressions at each spot, we estimate the abundance of all cell types in each spot ($W_{S \times C}$) and gene-specific scaling parameters ($M_{G \times 1}$) using Cell2location \citep{Kleshchevnikov2022} that takes as input the ST data $X^{st}$ and reference cell type signatures ($\mathcal{G}^{C}: C \times G$) estimated from the single-cell data $X^{sc}$. We then compute the absolute number of mRNA molecules contributed by each cell type $c$ to each spot $s$ to quantify the gene expression of a specific gene $g$ at spot $s$ for cell type $c$ as given by:
\begin{equation}
    U_{S \times C \times G} = \{u_{scg}\} = W*\mathcal{G}^{C}*M
\end{equation}
where $*$ represents matrix multiplication without sum-reduction. For single-cell resolution datasets, the cell type annotation of the spots are directly used as cell type abundance information, where $W_{S \times C}$ represents the cell type annotation of each spot and $U_{S \times C \times G}$ represents the original single-cell resolution ST dataset. This gives us,\\
\begin{equation}
    U_{S \times C \times G} = U_{S \times G} = \{u_{sg}\} = X^{st}
\end{equation}

% We feed $X^{st}$, $X^{sc}$ and its corresponding cell type annotations to cell2location \citep{Kleshchevnikov2022} to acquire a gene-specific scaling parameter ($M: G \times 1$), the absolute density of cell type $c$ in spot $s$ ($W: S \times C$) and reference cell type signatures ($\mathcal{G}^{C}: C \times G$). We then disentangle cell2locations formulation to calculate the absolute number of mRNA molecules contributed by each cell type $c$ to each spot $s$, to acquire the formulation for the inferred gene expression of a specific gene $g$ at spot $s$ for cell type $c$ as follows,\\
% $$U_{S \times C \times G} = \{u_{scg}\} = W*\mathcal{G}^{C}*M$$
% where $*$ represents matrix multiplication without sum-reduction.\\
% \\
The expected expression of a specific gene $g$ at spot $s$ for cell type $c$ can then be calculated as:
\begin{equation}
    E_{S \times C \times G} = \{e_{scg}\} = \frac{\Sigma_{n \in C} \Sigma_{s' \in S} u_{s'ng}}{S} \times \frac{\Sigma_{g' \in G} u_{scg'}}{\Sigma_{n \in C} \Sigma_{g' \in G} u_{sng'}}
\end{equation}
% $$$$


For single cell resolution datasets, since each cell is of a specific cell type the above formulation collapses to,\\
\begin{equation}
    E_{S \times C \times G} = E_{C \times G} = \{e_{cg}\} = \frac{\Sigma_{s' \in S} u_{s'cg}}{S}
\end{equation}


%E_{S \times C \times G} = \{e_{scg}\} = \frac{\Sigma_{n=1}^C \Sigma_{s'}^S u_{s'ng}}{S} \times \frac{\Sigma_{g'=1}^G u_{scg'}}{\Sigma_{n=1}^C \Sigma_{g'=1}^G u_{sng'}}

\subsection*{Inference of spatial communication domains based on ligand-target neighborhood scores}
In order to identify distinct spatial communication domains, which we define as a group of spots harboring similar ligand-target interactions, we performed clustering of the spots based on the ligand-target neighborhood activity scores. After computing the ligand-target neighborhood scores for all curated ligand-target pairs, each spot $s$ can be represented as a vector of ligand-target activity scores, $\CA_{l,t}$ of ligand-target pairs $(l,t) \in \CS_{lt}$ and based on these ligand-target activity score vectors, spots are clustered using the scanpy library \citep{Wolf2018}. First, the top 2000 highly variable ligand-target pairs are identified followed by principal component analysis. Leiden clustering is performed on the k-nearest neighbor graph obtained based on the principal components. 

% We clustered the spots with respect to the neighborhood scores generated to acquire domains with similar ligand-target interactions. The scanpy library is adopted to perform clustering over the top 2000 highly variable ligand-target pairs found using scanpys highly\_variable\_genes function. We perform principle component analysis over the highly variable genes using scanpys pca function. We then compute the neighborhood graph using scanpys neighbor function with the number of neighbors parameter set to 20. We then perform leiden clustering at dataset specific resolutions using scanpys leiden clustering function. The clusters have been generated for TNBC and Fetal Liver samples at a resolution of 0.6 and 0.3 respectively. The same has been generated for Mouse Brain at a resolution of 1.3 after integrating the samples as mentioned in section [reference].\\

% To exhibit the discreteness of each domain, we calculated the cosine similarity amongst spots in three different settings, (i) by using only the neighborhood scores for each spot, (ii) by using only the cell type abundance for each spot and (iii) by using a combination of both the neighborhood scores and cell type abundance for each spot. In particular, the cosine similarity between two spots given the neighborhood scores and cell type abundance across the two spots are calculated to be,\\
% $$k(\mathbf{x},\mathbf{y}) = \frac{\Sigma_{i \in N} x_i y_i}{\sqrt{\Sigma_{i \in N} x_i^2}\sqrt{\Sigma_{i \in N} x_i^2}}$$
% where $\mathbf{x}$ and $\mathbf{y}$ represents a vector of $N$ ligand-target scores or $C$ cell type abundances at the two spots respectively. The formulation produces a higher value if the cell type abundance/ligand-target neighborhood score between two spots are very similar and low otherwise.

\paragraph{Inferring cell type-specific sub-communication domains:} Additionally, we further sub-clustered the spots present within each of the spatial communication domains based on the cell type abundance of the spots. For sub-clustering also, Leiden clustering was adopted and implemented using scanpy.

\subsection*{Inference of communication domain-specific ligand-target activities}

\paragraph{Differential ligand-target activity across communication domains:} The Wilcoxon-rank sum test was utilized to estimate the ranked, differentially active ligand-target pairs for each communication domain using scanpy's `rank\_genes\_groups()' function.

\paragraph{Ranking of ligands based on activity within communication domains:} Renoir further allows the ranking of ligand activity within a communication domain. In order to perform ranking of ligand activity, we first determine the major cell types (cell types that together constitute $\geq 90\%$ of cell type abundance) within a communication domain based on the average proportion of the cell types across all the spots within that domain. For each major cell type in a domain, we identify the top marker genes ($n=500$) from the scRNA-seq dataset. We acquire a list of  ligand-receptor pairs from NATMI's connectomeDB \citep{Hou2020} and identify the major ligands that are also present within the marker gene lists of the major cell types within the domain. For each ligand, we identify a set of marker targets across the major cell types within the domain. A target is selected if it falls within the marker gene list of the cell type and the cell type expresses a receptor for that ligand. For each ligand, we compute the average neighborhood activity score for each of its marker targets (averaged across the spots expressing the associated receptor cell type). Each ligand is then scored based on the Pearson correlation between the average neighborhood scores of the ligand-target pairs of the cell type marker targets and the regulatory potential scores for each ligand-target pair acquired from NicheNet's prior model \citep{Browaeys2020}. Using the ligand score, Renoir ranks the ligands, where the ranking represents the amount of agreement between the ligand-target activity in our dataset and the prior knowledge of the pairs. While in our analyses, we used ligand-receptor and ligand-target pairs from specific sources, Renoir can also work with user-defined ligand-receptor and ligand-target pairs.
% every expressed receptor and its corresponding cell type, a list of marker targets (targets which are markers of the receptor cell type) are identified for each ligand associated with the receptor, giving a list of ligand-receptor associated ligand-target pairs. 
% A simple average of the neighborhood scores across the spots which contain the cell types that the ligand-receptor associated ligands and targets are markers of, is found for each ligand-target pair to acquire an average ligand-target neighborhood score across spots expressing the associated receptor cell type. With the list of ligand-target pairs and their associated average neighborhood scores, we then score each ligand by the correlation between the average neighborhood scores of the ligand-target pairs and the regulatory potential scores for each ligand-target pair acquired from NicheNets prior model \citep{Browaeys2020}. The ligand score then allows Renoir to rank the ligand-target pairs, where the ranking implies how in tune are the datasets ligand-target interactions with our prior knowledge on the pairs. While the analysis in this work used ligand-receptor pairs and ligand-target pairs from specific sources, Renoir will also work with user defined ligand-receptor and ligand-target pairs.

\paragraph{Inferring major cell type-specific ligand-target activity within a communication domain:} To identify cell type-specific ligand-target activity within each communication domain, we first considered the ligand-target pairs for which both the ligand and target fell within the top 200 markers for the major cell types of the communication domain. For each domain where the ligand and target were markers of the major cell types, a simple average of the Pearson correlations between the neighborhood activity scores of the ligand-target pair and the cell type proportions of the cell types expressing the ligand and target respectively were calculated across the spots pertaining to the domain being considered. The correlation score measures how prominent a cell type-specific ligand-target activity is within a communication domain.

% the top 100 marker genes of each cell type using the available scRNA-seq data. All ligand-target pairs for which both the ligand and target fell within the top 100 markers were considered. A threshold $n$ is then considered after which the top $n$ communication domains were assigned to each cell type on the basis of average cell type proportions across each cluster. A threshold of 3 was considered for the datasets utilized in this work. For each domain where the ligand and target were markers of the cell types the domain was prominent with, a simple average of the pearson correlations between the neighborhood scores of the ligand-target pair and the cell type proportions of the cell types the ligand and target were markers of were calculated across the spots pertaining to the domain being considered. This gives a measure of how prominent a cell type specific ligand-target activity is within a communication domain.

\subsection*{Inference of spatial mapping of pathway activity}

\paragraph{Construction of genesets of known pathways or Denovo clusters:} Renoir further allows the spatial mapping of a pathway activity by utilizing existing or denovo gene sets. First, for a particular pathway represented by a gene set, we determine the ligand-target pairs active in the dataset by intersecting the set of all possible gene pairs within a gene set with the set of curated ligand-target pairs. In this work, we restricted our usage to the GSEA Molecular Signatures Database \citep{doi:10.1073/pnas.0506580102, 10.1093/bioinformatics/btr260}, specifically HALLMARK and CANONICAL PATHWAYS (KEGG, WIKIPATHWAYS) gene sets. Renoir can also utilize any predefined gene set provided by the user, specific to the user's requirements. Apart from utilizing gene sets for existing pathways, Renoir can also identify de novo clusters of ligand-target pairs based on their similarity of the neighborhood activity scores across the spots. For performing the clustering of ligand-target pairs, Renoir employs hierarchical clustering strategies, namely HDBSCAN and dynamic hclust. HDBSCAN was implemented using the hdbscan python package with the min\_pts parameter set to the lowest possible value and dynamic hclust was implemented using the dynamicTreeCut package.

% We define a number of genesets with respect to existing pathways using the following procedure. Each gene within a gene set is paired with each other to produce hypothetical ligand-target pairs. The hypothetical pairs for a specific gene set and the existing ligand-target pairs present within the neighborhood scores are intersected to gain a set of ligand-target pairs specific to the gene set and dataset. In this work, we restricted our usage to the GSEA Molecular Signatures Database, specifically HALLMARK and CANONICAL PATHWAYS (KEGG, WIKIPATHWAYS) gene sets. Renoir can also utilize any predefined geneset provided by the user, specific to the users requirements. De Novo clusters are generated by clustering ligand-target interactions present within the spot-wise min-max normalized neighborhood scores. To do so, one of two clustering approaches are considered, namely HDBSCAN and dynamic hclust. HDBSCAN was implemented using the hdbscan python package with the min\_pts parameter set to the lowest possible value (see section [reference]) and dynamic hclust was implemented using the dynamicTreeCut package. Every other parameter followed default settings. The key difference seen amongst the two approaches across a specific dataset is the size and structure of clusters generated. HDBSCAN produces large number of clusters with small number of ligand/target specific interactions and dynamic hclust produces relatively low number of clusters with larger ligand/target specific interactions, thus providing different perspectives to ligand/target specific interactions across the spatial topology.

\paragraph{Spatial mapping of existing pathways and Denovo clusters:}
To generate a spatial map of previously known pathways or Denovo clusters, we computed the first principal component after performing dimensionality reduction over the neighborhood activity scores of the subset of ligand-target pairs corresponding to a pathway/Denovo cluster using the `TruncatedSVD' function of sklearn. This generates a pathway activity score for each spot in the spatial transcriptomics data. 
% Additionally, the De Novo clusters are further refined to consider only those ligand-target pairs that substantially affect the overall principle component score. To do so, we first calculate the knee point over the first (and only) right singular vector of the De Novo cluster. The ligand-target pairs are then threshold against the knee point.

%Every downstream task further discussed have been performed over all datasets investigated. Spotwise-cell type specific clusters (section [reference]) have been generated for the Breast Cancer and Fetal Liver samples at a resolution of 0.5 and 0.6 respectively. No hyperparameters were set for the Pairwise-gene specific De Novo clusters (section [reference]) via hierarchical clustering with dynamic tree cut (dynamic hclust). A minimum resolution of 0.3 is considered by default while using the HDBSCAN approach. Each De Novo cluster and pathway visualised across the spatial topography were considered only if each cluster or pathway had a minimum of 5 or 6 ligand-target pairs respectively.

%\subsection{Overview of ligand-target activity across a dataset}
%To generate an overview of multiple ligand-target activity across the dataset, a Sankey plot was developed to infer the relationship between cell types, ligand-target pairs and communication domains. The diagram consists of five levels, which initiates with the relationship between the distribution of cell types (level 1) across every communication domain (level 2). The nodes within level 1 represents unique cell types and the nodes within level 2 represents communication domains. Each edge from level 1 to level 2 represents the average proportion of each cell type within a communication domain. Given the unique set of cell types $CT$ and the abundance value of each cell type across the spots $S$, the average proportion of cell type $ct$ at domain $R$ is calculated as so,
%$$avg(ct,R) = \frac{\Sigma_{s \in S_R} p_{ct}(s)}{|S_R|}$$
%$$p_{ct}(s) = \frac{w_{s,ct}}{\Sigma_{c \in CT} w_{s,c}}$$
%The nodes within level 3 represents ligands and each edge between level 2 and level 3 represents the cumulative sum of the average neighborhood scores of a specific ligand (with varying targets) across each communication domain defined by,
%$$Edge(l,R) = \Sigma_{t} avg(lt,R)$$
%$$avg(lt,R) = \frac{\Sigma_{s \in R} NeighbScore_{l,t}(s)}{|S_R|}$$
%Nodes within level 4 represent targets with each edge from level 3 to level 4 representing the cumulative sum of the average ligand-target neighborhood scores across each communication domain defined by,
%$$Edge(l,t) = \Sigma_{R} avg(lt,R)$$
%Nodes within level 5 represent communication domains once again and each edge from level 4 to level 5 represents the cumulative sum of the average neighborhood scores of a specific target (with varying ligands) across each communication domain, defined by
%$$Edge(R,t) = \Sigma_{l} avg(lt,R)$$

\subsection*{Benchmarking of Renoir}

\subsubsection*{Generation of simulated data}
For benchmarking the performance of Renoir, we simulated spatial datasets based on a semi-synthetic approach inspired from the simulation approach in \citep{Liu2022}. In order to simulate a dataset with known spatial interaction between ligand and target, we start with a paired scRNA-seq and ST data (the reference dataset) and cell type deconvolution of the reference ST dataset is performed using cell2location \citep{Kleshchevnikov2022}. After that, we simulate the ST dataset using the following steps in order to maintain a controlled environment, where we can keep track of the different kinds of interactions enriched within the data. 
% whilst ensuring to the best extent that there is no noise to affect the expected results.\\
\begin{enumerate}
    \item Select a set of ligand-target pairs $LT$ such that both the ligand and target are marker genes (top 100 markers) of at least one cell type in the paired scRNA-seq data. After that, the pairs in $LT$ are segregated into two mutually exclusive subsets: $LT_{m}$ and $LT_{s}$. $LT_{m}$ consists of the ligand-target pairs where the target gene is a marker for multiple (two or more) cell types and $LT_{s}$ consists of every other pair, where the target gene is a marker of exactly one cell type. In both $LT_{m}$ and $LT_{s}$, the ligand could be a marker of one or more cell types. In our analysis, we were able to find a maximum of two cell types associated with each of the ligand or target genes.
    \item For a specific ligand-target pair $lt \in LT_{s}$, where the target is a marker of a cell type $ct_{t1}$; in the ST dataset, we generate two types of reference spots, positive and negative spots respectively. Positive reference spots denote the spots within the tissue where ligand-target interactions for $lt$ are simulated through the enrichment of the ligand cell type and target cell type $ct_{t1}$, where a spot is enriched with a cell type by adding $n_e = Uniform(2,5)$ single cells of the said cell type into said spot. The UMI counts are then scaled back to the original UMI count to maintain the original UMI distribution of the sample being simulated. Negative references refer to the spots within the tissue that do not harbor any activity for $lt$ as indicated by the low proportion of ligand cell type and target cell type $ct_{t1}$. Specifically, negative reference spots are selected as a set of spots for which both ligand cell type and target cell type $ct_{t1}$ proportions are low and their neighbors also have low proportions (below acceptable threshold) of ligand and target cell types. In the negative reference spots, interactions of $lt$ can not occur given the low abundance of both the ligand and target cell types in these spots as well as in the neighboring spots.
    % This helps provide a set of true negative like samples, where we assume that the ligand-target interactions pertinent to the ligand and target cell types could not occur since the ligand and target cell types are not abundant across the negative samples. The number of spots selected varies across datasets.
    \item For a specific ligand-target pair $lt \in LT_{m}$ where the target is a marker of two cell types $ct_{t1}$ and $ct_{t2}$; in the corresponding ST dataset, we generate two types of reference spots, positive and negative. Without loss of generality, we assume that the cell type $ct_{t1}$ does not express a known receptor of the ligand $l$ but the cell type $ct_{t2}$ expresses a known receptor of $l$ (according to a ligand-receptor database, e.g., NATMI). Positive references are those spots within the tissue that have been enriched with the ligand cell type and target cell type $ct_{t2}$. Specifically, a set of spots are selected that are not devoid of the ligand cell type. This is to ensure that the enrichment follows the underlying distribution of cell types in the dataset. For each of the selected spots, the spot is enriched with the ligand cell type and the neighbors of the spot are enriched with the target cell type $ct_{t2}$. Negative references are those spots within the tissue that have been enriched with the ligand cell type and target cell type $ct_{t1}$ (type 1 negative); or spots within the tissue that do not harbor any activity for $lt$ as indicated by the low proportion of ligand cell type and target cell type $ct_{t2}$ (type 2 negative). To simulate type 1 negative spots, a set of spots are selected from the set of all spots with low $ct_{t1}$ and $ct_{t2}$ proportions. We then enrich these spots with the ligand cell type and the target cell type $ct_{t1}$. Thus for the type 1 negative spots the ligand-target interactions pertinent to the ligand cell type and target cell type $ct_{t1}$ could not occur since no known receptors are expressed by $ct_{t1}$ despite the co-localization of the cell types. The simulation of type 2 negative spots follows the same strategy as in case of $lt \in LT_{s}$ considering $ct_{t2}$ as the target cell type. 
    % \textcolor{red}{Is this strategy correct?? For a target involving more than 1 cell type, we have 2 types of negative spots - regular negative i.e., low abundance, special negative: where cell type is there but does not contain a receptor. Rewrite this section accordingly. Also, how are you enriching a cell type? That is not written here.}
    
\end{enumerate}

\subsubsection*{Scoring and evaluation strategy}
To evaluate the performance of a CCC method in inferring the ligand-target activity at spot resolution, we employ the following scores. 
\paragraph{Accuracy of spatial activity:} The true activity of a ligand-target pair across the positive and negative reference spots is computed by employing the local Moran's bivariate test-statistic ($R_{s}^{ct_l ct_t}$) between the ligand and target cell types,\\
$$R_{s}^{ct_l ct_t} = \left(\frac{w_{s,ct_l} - \mu_{ct_l}}{\sigma_{w_{s,ct_l}}} \times SL_{s, ct_t} \right) + \left(\frac{w_{s,ct_t} - \mu_{ct_t}}{\sigma_{ct_t}} \times SL_{s, ct_l} \right)$$

$$SL_{s, x} = \frac{\sum_{i \in \CN_s} x_i - \mu_x \cdot |\CN_s|}{\sigma_x} $$

where $w_{s,x}$ represents the cell type abundance of cell type $x$ at spot $s$, $\mu_x$ represents the mean cell type abundance of cell type $x$ across all spots considered, $ \sigma_x$ represents the standard deviation of cell type abundance of cell type $x$ across all spots considered and $SL_{s, x}$ represents the spatial lag of cell type $x$ at spot $s$. This provides a measure of local interactions between the ligand and target cell types. Pearson's correlation (as discussed before) is then used to calculate the average correlation between the ligand-target activity scores inferred by each method and $R_{s}^{ct_l ct_t}$ which gives a measure of how well does the inferred activity scores align with the true interactions as quantified in terms of spatial cell type co-localization.

\paragraph{Precision and Recall:} To calculate the precision and recall of a CCC inference method in inferring the spatial interactions of ligand and target, we divide the spots into two sets based on a predefined threshold set on the ligand-target scores computed by the CCC inference method. Given a specific method, the spots which have a score greater than or equal to the threshold are considered to be positive predicted spots (harboring the ligand-target interaction) and the spots with a score lesser than the threshold are considered to be negative predicted spots (no ligand-target interaction). We then consider the these positive and negative predicted spots and compare them against the ground truth positive and negative reference spots. A predicted spot is considered to be a true positive (TP) if the spot is a positive predicted spot and is also a positive reference spot, true negative (TN) if the spot is a negative predicted spot and is also a negative reference spot, false positive (FP) if the spot is a positive predicted spot but is a negative reference spot and a false negative (FN) if the spot is a negative predicted spot but a positive reference spot. For each method, we calculate the TP, TN, FP and FN values to calculate the precision and recall of a method for a given simulated dataset. Thus precision and recall are calculated for each simulated ligand-target pair and the final precision and recall for a dataset for a specific threshold is computed to be the average of the precision and recall values across all considered ligand-target pairs for the same dataset. Instead of selecting a single threshold, the threshold is varied to determine the precision-recall curve for each method across various threshold values. 

\paragraph{Comparison of communication domains with ground truth layers in DLPFC dataset}
In order to estimate how coherent the inferred communication domains are with known regions of the dataset, we first estimate the communication domains across all 12 DLPFC samples for each of the methods (Renoir, COMMOT, SpatialDM and stLearn). This is done so by applying leiden clustering (resolutions range from 0.1 to 1.6 in intervals of 0.1) over the ligand-target interaction scores generated by each of the methods across all samples. The NMI and ARI values are calculated between the cluster labels and ground truth annotations; for each of the methods across all samples over all resolutions. At each resolution, the averaged NMI and ARI values are calculated for each method across all samples. For each method, the cluster labels corresponding to the resolution at which the averaged NMI and ARI values is the best for said method in each sample, is now considered as the best communication domain inferred by a method for a specific sample. This results with a list of best NMI/ARI values for each method across each sample. The box plots are then generated over the final inferred NMI/ARI values.

\subsection*{Preparation of fetal liver data}
FFPE fetal liver tissue was obtained under ethical approval from the Centralised Institutional Research Board of the Singapore Health Services, Singapore SingHealth and National Health Care Group Research Ethics Committees. To determine the quality of RNA in FFPE tissue, RNA was extracted from 10$\mu m$ sections using a Qiagen FFPE RNA extraction kit (\#73504) and following the manufacturer's protocol. RNA quality was determined as ``good" if the percentage of RNA fragments $>$200 nucleotides (DV200\%) was at least 50\% of all fragments. Next, a 5$\mu m$ tissue section was placed on spatially barcode oligonucleotides, within the capture area (6.5 x 6.5$mm$) of the Visium slide (4 capture areas/slide) and processed according to the manufacturer's protocol. Briefly, Visium slides with affixed FFPE tissue sections were deparaffinised in fresh xylene and rehydrated in reducing ethanol (in water) concentration gradient. The slides were stained with Haematoxylin and Alcoholic Eosin and imaged at 10 times optical magnification using a Nikon Ni-E Microscope. RNA was then de-crosslinked and tissue sections were hybridised with Visium Human Transcriptome Probe Set, overnight at 50\textdegree C. The following day, the slides were washed with supplied post-hybridisation wash buffer and hybridized probe pairs were ligated using supplied ligase buffer at 37\textdegree C. After ligation, the slides were washed with room temperature and pre-heated post-ligation wash buffer. The ligated probes were then released by digesting hybridized RNA with supplied RNase enzyme at 37\textdegree C. The tissue sections were then permeabilised at 37\textdegree C using the supplied permeabilising enzyme. Permeabilization enabled ligated probes to be captured by spatially barcoded oligonucleotides immobilised on the Visium slide. Captured hybridised probes were then extended using probe extension mix at 45\textdegree C and extended probes were eluted accordingly. To determine appropriate amplification cycle numbers for library preparation in the next step, 1$\mu l$ of eluted probes were quantitively amplified (qPCR) in accordance with the manufacturer's instructions. Quantification cycle (Cq) value at ~25\% peak relative fluorescence unit (RFU) was used to set cycle numbers during library preparation. The library was then quantified using TapeStation (Agilent High Sensitivity D5000) to determine the average fragment size and concentration before sequencing.

\subsection*{Integration of mouse brain datasets}
The two mouse brain samples were integrated using the stLearn workflow \citep{Pham2020.05.31.125658}. Neighborhood scores generated for each mouse brain sample were concatenated and min-max normalized across ligand-target pairs for every spot. The principal components calculated over the spatial transcriptomic data of each sample via scanpy were used as input to harmony algorithm \citep{Korsunsky2019}, which projected the spots onto a shared embedding space and thus presenting a new set of adjusted principal components for the concatenated dataset. Leiden clustering was performed via scanpy over the concatenated dataset with adjusted principal components at a resolution of 1.3 to obtain a set of spatial communication domains shared across the two brain samples.

% \subsubsection{Frequency table: domain and region labels}
% We performed cross tabulation between the spots annotated with the spot-wise interaction cluster labels and the spots annotated with available spot-wise domain annotations for the two Mouse brain samples, to acquire a frequency table denoting the number of spots per domain vs the number of spots per region. Cross tabulation was performed using pandas crosstab function, which was passed onto seaborn clustermap function to generate the final cluster map of the frequency table. 

\subsection*{Mapping of regional astrocyte subtype activity in mouse brain}
In order to find localized astrocyte subtype activity in appropriate regions of the mouse brain dataset, we first identified the top 100 marker genes for every cell type (including astrocyte subtypes) using the available mouse brain snRNA-seq dataset. We then identified the top 100 differentially active ligand-target pairs for each region of the Mouse Brain dataset. Two sets of pairs were then generated, (i) the first set consisted of ligand-target pairs where the ligand was a marker gene of a specific astrocyte subtype, the ligand-target pair displayed differential activity in a region associated with said astrocyte subtype and the target was a marker of a cell type associated with said region (ii) the second set consisted of ligand-target pairs where the target was a marker gene of a specific astrocyte subtype, the ligand-target pair was differentially active in a region associated with said astrocyte subtype and the ligand was a marker of another cell type associated with said region. The above steps were carried out using the scanpy library `rank\_genes\_group' function.

\subsection*{Correlation measures used for analyzing \textit{PLG:MARCO} interaction in fetal liver}

The following correlation measures were utilized to measure the similarity between \textit{PLG:MARCO} neighborhood scores with \textit{PLG}$^{+}$ and \textit{PLG}$^{-}$ hepatocytes and \textit{FOLR2}$^{+}$ Macrophages for the Fetal Liver dataset,
\begin{itemize}
    \item \emph{Pearson correlation}: We utilized Pearson correlation to find the correlation amongst specific ligand/target expression values (either derived from scRNA-seq data or spatial transcriptomic data) and cell type abundances. We also used Pearson correlation to acquire an initial estimate of the correlation amongst the neighborhood scores of a specific ligand-target pair with the abundances of the cell types expressing the ligand and the target. The overall correlation measure can be summarised as,\\
    $$r_{{l,t},{ct_l,ct_t}} = \frac{P_{lt}(ct_l) + P_{lt}(ct_t)}{2}$$\\
    where,\\
    \\
    $$P_{lt}(ct_x) = \frac{\Sigma_{i \in S_x} (NS_{lt}(i) - \mu_{NS_{lt}})(w_{i,ct_x} - \mu_{w_{ct_x}})}{\sqrt{\Sigma_{i \in S_x} (NS_{lt}(i) - \mu_{NS_{lt}})^2 \Sigma_{i \in S_x} (w_{i,ct_x} - \mu_{w_{ct_x}})^2}}$$
    \\
    $l$ and $t$ represent a specific ligand and target respectively\\
    $ct_l$ and $ct_t$ represent the cell types expressing the ligand and the target respectively (as a marker) \\
    $NS_{lt}(s)$ represents the neighborhood scores of ligand-target pair $lt$ at spot $s$\\ 
    $NS_{lt}$ represents a vector of neighborhood scores of ligand-target pair $lt$ across all $S_x$ spots\\
    $w_{s,y}$ represents the cell type abundance of the cell type the ligand($ct_l$) or target($ct_t$) is a marker of at spot $s$\\ 
    $w_{ct_x}$ is a vector of cell type abundances of the cell type the ligand($ct_l$) or target($ct_t$) is a marker of, across all $S_x$ spots\\
    $\mu_{NS_{lt}}$ is the average neighborhood score for a specific ligand-target pair $lt$\\
    $\mu_{w_{ct_x}}$ is the average cell type abundance for a specific cell type $ct_x$\\
    $S_x = \{s_j, (NS_{lt}(s_j) > 0) | (w_{s_j,ct_x} > 0)\}$ where $x$ is either a ligand $l$ or target $t$.\\
    \item \emph{Pearson correlation over spatial lag}: We furher used Pearson correlation over spatial lag as a measure of spatial correlation \citep{Lee2001}. Pearson correlation over spatial lag considers the sum of expression or abundance values over a spot and its corresponding neighbors before applying Pearson correlation. Formally, Pearson correlation in the context of expression or abundance values can be defined as,\\
    $$r_{xy} = \frac{\Sigma_{i \in S} (X_i - \mu_{X})(Y_i - \mu_{Y})}{\sqrt{\Sigma_{i \in S} (X_i - \mu_{X})^2 \Sigma_{i \in S} (Y_i - \mu_{Y})^2}}$$
    $X$ and $Y$ are the spatial lags defined by,\\
    $$X_i = \Sigma_{n \in N(x_i)} x_n$$\\
    where $N(x_i)$ represents the neighborhood of $x_i$.
    $X = {x_i}, i \in S$; $S_x = \{s_j, (X_{s_j} > 0) | (Y_{s_j} > 0)\}$
    The spatial lags are calculated with respect to gene expressions from the spatial transcriptomics data or the spatial map of cell type abundances ($x$ or $y$). While considering neighborhood scores, the scores are considered as they are since the neighborhood score at each spot takes into account the neighbors of the spot.
    \item \emph{Lee's L statistic}: We utilised Lee's L statistic bivariate spatial association measure to incorporate spatial autocorrelation alongside spatial patterning between the two variables of interest \citep{Lee2001}. In our case, we adopted the measure to view the spatial co-patterning of neighborhood scores, gene expressions from the spatial transcriptomics data or the spatial map of cell type abundances with one another. This is defined as,
    $$L_{X,Y} = \sqrt{L_{X,X} L_{Y,Y}} r_{xy}$$
    where,\\
    $$L_{X,X} = \frac{\Sigma_{i \in S} (X_i - \mu_x)^2}{\Sigma_{i \in S} (x_i - \mu_x)^2}$$
\end{itemize}